\section{Procedures e Funcoes}

\subsection{Funcoes de Negocio}

\subsubsection{Calculo Splitwise}

\begin{lstlisting}[caption=Funcao para calcular divisao Splitwise]
CREATE OR REPLACE FUNCTION calcular_splitwise_grupo(
    p_grupo_id UUID,
    p_data_inicio DATE DEFAULT DATE_TRUNC('month', CURRENT_DATE)::DATE,
    p_data_fim DATE DEFAULT (DATE_TRUNC('month', CURRENT_DATE) + INTERVAL '1 month - 1 day')::DATE
)
RETURNS TABLE (
    usuario_id UUID,
    nome_usuario VARCHAR(100),
    total_pago DECIMAL(15,2),
    valor_devido DECIMAL(15,2),
    saldo DECIMAL(15,2)
) AS $$
DECLARE
    total_gastos DECIMAL(15,2);
    num_usuarios INTEGER;
    valor_por_pessoa DECIMAL(15,2);
BEGIN
    -- Verificar se o grupo usa modelo Splitwise
    IF NOT EXISTS (
        SELECT 1 FROM grupos 
        WHERE id = p_grupo_id AND modelo_organizacao = 'splitwise'
    ) THEN
        RAISE EXCEPTION 'Grupo nao configurado para modelo Splitwise';
    END IF;
    
    -- Contar usuarios ativos no grupo
    SELECT COUNT(*) INTO num_usuarios
    FROM usuario_grupos ug
    WHERE ug.grupo_id = p_grupo_id AND ug.ativo = true;
    
    -- Calcular total de gastos do periodo
    SELECT COALESCE(SUM(t.valor), 0) INTO total_gastos
    FROM transacoes t
    WHERE t.grupo_id = p_grupo_id
      AND t.tipo = 'despesa'
      AND t.status = 'confirmada'
      AND t.data_transacao BETWEEN p_data_inicio AND p_data_fim;
    
    -- Calcular valor devido por pessoa
    valor_por_pessoa := total_gastos / num_usuarios;
    
    -- Retornar resultado
    RETURN QUERY
    SELECT 
        u.id,
        u.nome,
        COALESCE(SUM(t.valor), 0) AS total_pago,
        valor_por_pessoa AS valor_devido,
        (valor_por_pessoa - COALESCE(SUM(t.valor), 0)) AS saldo
    FROM usuarios u
    JOIN usuario_grupos ug ON u.id = ug.usuario_id
    LEFT JOIN transacoes t ON u.id = t.usuario_id 
        AND t.grupo_id = p_grupo_id
        AND t.tipo = 'despesa'
        AND t.status = 'confirmada'
        AND t.data_transacao BETWEEN p_data_inicio AND p_data_fim
    WHERE ug.grupo_id = p_grupo_id AND ug.ativo = true
    GROUP BY u.id, u.nome, valor_por_pessoa
    ORDER BY u.nome;
END;
$$ LANGUAGE plpgsql;
\end{lstlisting}

\subsubsection{Atualizacao de Saldos}

\begin{lstlisting}[caption=Procedure para atualizar saldos de todas as contas]
CREATE OR REPLACE PROCEDURE atualizar_saldos_usuario(p_usuario_id UUID)
LANGUAGE plpgsql AS $$
DECLARE
    conta_record RECORD;
    novo_saldo DECIMAL(15,2);
BEGIN
    -- Percorrer todas as contas do usuario
    FOR conta_record IN 
        SELECT id, saldo_inicial, tipo
        FROM contas 
        WHERE usuario_id = p_usuario_id AND ativa = true
    LOOP
        -- Calcular novo saldo
        SELECT conta_record.saldo_inicial + COALESCE(SUM(
            CASE 
                WHEN t.tipo = 'receita' THEN t.valor
                WHEN t.tipo = 'despesa' THEN -t.valor
                WHEN t.tipo = 'transferencia' AND t.conta_id = conta_record.id THEN -t.valor
                WHEN t.tipo = 'transferencia' AND t.conta_destino_id = conta_record.id THEN t.valor
                ELSE 0
            END
        ), 0) INTO novo_saldo
        FROM transacoes t
        WHERE (t.conta_id = conta_record.id OR t.conta_destino_id = conta_record.id)
          AND t.status = 'confirmada';
        
        -- Atualizar saldo
        UPDATE contas 
        SET saldo_atual = novo_saldo,
            updated_at = CURRENT_TIMESTAMP
        WHERE id = conta_record.id;
        
        -- Log da atualizacao
        INSERT INTO audit.audit_logs (tabela, operacao, usuario_id, dados_novos)
        VALUES ('contas', 'UPDATE', p_usuario_id, 
                jsonb_build_object('conta_id', conta_record.id, 'novo_saldo', novo_saldo));
    END LOOP;
    
    COMMIT;
END;
$$;
\end{lstlisting}

\subsubsection{Processamento de Recorrencias}

\begin{lstlisting}[caption=Funcao para processar transacoes recorrentes]
CREATE OR REPLACE FUNCTION processar_recorrencias()
RETURNS INTEGER AS $$
DECLARE
    rec_record RECORD;
    nova_transacao_id UUID;
    contador INTEGER := 0;
    proxima_data DATE;
BEGIN
    -- Processar recorrencias vencidas
    FOR rec_record IN 
        SELECT * FROM recorrencias 
        WHERE ativa = true 
          AND proxima_execucao <= CURRENT_DATE
          AND (data_fim IS NULL OR proxima_execucao <= data_fim)
    LOOP
        -- Criar nova transacao
        INSERT INTO transacoes (
            descricao, valor, tipo, data_transacao,
            conta_id, categoria_id, usuario_id, grupo_id,
            status, observacoes
        ) VALUES (
            rec_record.nome,
            rec_record.valor,
            rec_record.tipo,
            rec_record.proxima_execucao,
            rec_record.conta_id,
            rec_record.categoria_id,
            rec_record.usuario_id,
            rec_record.grupo_id,
            CASE WHEN rec_record.auto_confirmar THEN 'confirmada' ELSE 'pendente' END,
            'Transacao recorrente gerada automaticamente'
        ) RETURNING id INTO nova_transacao_id;
        
        -- Calcular proxima execucao
        SELECT CASE rec_record.frequencia
            WHEN 'diario' THEN rec_record.proxima_execucao + INTERVAL '1 day'
            WHEN 'semanal' THEN rec_record.proxima_execucao + INTERVAL '1 week'
            WHEN 'quinzenal' THEN rec_record.proxima_execucao + INTERVAL '15 days'
            WHEN 'mensal' THEN rec_record.proxima_execucao + INTERVAL '1 month'
            WHEN 'bimestral' THEN rec_record.proxima_execucao + INTERVAL '2 months'
            WHEN 'trimestral' THEN rec_record.proxima_execucao + INTERVAL '3 months'
            WHEN 'semestral' THEN rec_record.proxima_execucao + INTERVAL '6 months'
            WHEN 'anual' THEN rec_record.proxima_execucao + INTERVAL '1 year'
            ELSE rec_record.proxima_execucao + INTERVAL '1 month'
        END INTO proxima_data;
        
        -- Atualizar recorrencia
        UPDATE recorrencias 
        SET proxima_execucao = proxima_data,
            updated_at = CURRENT_TIMESTAMP
        WHERE id = rec_record.id;
        
        contador := contador + 1;
        
        -- Log da operacao
        INSERT INTO audit.audit_logs (tabela, operacao, usuario_id, dados_novos)
        VALUES ('recorrencias', 'PROCESS', rec_record.usuario_id,
                jsonb_build_object('recorrencia_id', rec_record.id, 
                                 'transacao_id', nova_transacao_id,
                                 'proxima_data', proxima_data));
    END LOOP;
    
    RETURN contador;
END;
$$ LANGUAGE plpgsql;
\end{lstlisting}

\subsection{Funcoes de Relatorio}

\subsubsection{Resumo Financeiro}

\begin{lstlisting}[caption=Funcao para gerar resumo financeiro]
CREATE OR REPLACE FUNCTION gerar_resumo_financeiro(
    p_usuario_id UUID,
    p_data_inicio DATE,
    p_data_fim DATE
)
RETURNS TABLE (
    periodo TEXT,
    total_receitas DECIMAL(15,2),
    total_despesas DECIMAL(15,2),
    saldo_periodo DECIMAL(15,2),
    categoria_maior_gasto VARCHAR(100),
    valor_maior_gasto DECIMAL(15,2)
) AS $$
BEGIN
    RETURN QUERY
    WITH resumo_base AS (
        SELECT 
            TO_CHAR(t.data_transacao, 'YYYY-MM') as mes,
            SUM(CASE WHEN t.tipo = 'receita' THEN t.valor ELSE 0 END) as receitas,
            SUM(CASE WHEN t.tipo = 'despesa' THEN t.valor ELSE 0 END) as despesas
        FROM transacoes t
        WHERE t.usuario_id = p_usuario_id
          AND t.status = 'confirmada'
          AND t.data_transacao BETWEEN p_data_inicio AND p_data_fim
        GROUP BY TO_CHAR(t.data_transacao, 'YYYY-MM')
    ),
    maior_categoria AS (
        SELECT DISTINCT ON (TO_CHAR(t.data_transacao, 'YYYY-MM'))
            TO_CHAR(t.data_transacao, 'YYYY-MM') as mes,
            c.nome as categoria,
            SUM(t.valor) as total_categoria
        FROM transacoes t
        JOIN categorias c ON t.categoria_id = c.id
        WHERE t.usuario_id = p_usuario_id
          AND t.tipo = 'despesa'
          AND t.status = 'confirmada'
          AND t.data_transacao BETWEEN p_data_inicio AND p_data_fim
        GROUP BY TO_CHAR(t.data_transacao, 'YYYY-MM'), c.nome
        ORDER BY TO_CHAR(t.data_transacao, 'YYYY-MM'), SUM(t.valor) DESC
    )
    SELECT 
        rb.mes,
        rb.receitas,
        rb.despesas,
        (rb.receitas - rb.despesas),
        mc.categoria,
        mc.total_categoria
    FROM resumo_base rb
    LEFT JOIN maior_categoria mc ON rb.mes = mc.mes
    ORDER BY rb.mes;
END;
$$ LANGUAGE plpgsql;
\end{lstlisting}

\subsection{Funcoes de Validacao}

\subsubsection{Validacao de Limite de Credito}

\begin{lstlisting}[caption=Funcao para validar limite de credito]
CREATE OR REPLACE FUNCTION validar_limite_credito()
RETURNS TRIGGER AS $$
DECLARE
    saldo_atual DECIMAL(15,2);
    limite DECIMAL(15,2);
    tipo_conta VARCHAR(30);
BEGIN
    -- Buscar informacoes da conta
    SELECT c.saldo_atual, c.limite_credito, c.tipo
    INTO saldo_atual, limite, tipo_conta
    FROM contas c
    WHERE c.id = NEW.conta_id;
    
    -- Verificar apenas para cartao de credito
    IF tipo_conta = 'cartao_credito' AND NEW.tipo = 'despesa' THEN
        -- Calcular novo saldo apos a transacao
        IF (saldo_atual - NEW.valor) < -COALESCE(limite, 0) THEN
            RAISE EXCEPTION 'Transacao excede limite de credito. Limite: %, Saldo atual: %, Valor transacao: %',
                COALESCE(limite, 0), saldo_atual, NEW.valor;
        END IF;
    END IF;
    
    RETURN NEW;
END;
$$ LANGUAGE plpgsql;
\end{lstlisting}
